% 原书分章 PDF 第 21—24 页；第 20.8—20.9 节及致谢。
\section{小结}

本章为有限总体中因果相连个体所生成的纵向数据结构，建立了一般因果模型。随后，定义了由针对个体处理节点的干预所索引的反事实，以及相应的因果对比。在适当假设下，我们确立了：观察所有个体，或观察个体的随机样本时，这些因果量可从个体的观测数据中识别。我们的因果假设意味着，给定所有个体的过去时，时间 $t$ 的各个体条件独立。这使观测数据的似然可以分解，尽管数据由单次实验产生，并非来自独立实验的重复。为应对维数灾难，我们假设，一个个体对其他个体过去的依赖，可用有限维度量概括，且这种依赖在各个体间由同一个函数描述。随后介绍了数据分布的统计模型、统计目标参数，以及由此得到的统计估计问题。我们说明，可以使用交叉验证和超级学习，估计似然的不同因子。给定统计模型，以及识别某种干预下反事实均值的统计目标参数，推导了目标参数的有效影响曲线。我们表明，EIC 刻画极大似然估计量的正态极限分布，因此代表目标参数各估计量中的最优渐近方差。但是，由于维度较高，极大似然估计量在有限样本下无法良好定义，必须进行平滑。

平滑／正则化极大似然估计量未经靶向处理，因此对目标参数而言偏差过大。由此，它们通常无法得到服从渐近正态分布的统计目标参数估计量。因此，我们为该待估量构造了靶向极大似然估计量（targeted maximum likelihood estimators）。所提出 TMLE 的偏差是一个二阶项，涉及两个干扰参数的平方差 $\bar h_N-\bar h_0$ 和 $q_N-q_0$；这两个干扰参数为 $\bar h_0=\bar h(\bar g_0,\bar q_0)$ 及似然的相关因子 $\bar q_0$。

总的来说，我们认为，对这些个体动态网络因果模型的统计研究，是一个引人入胜且重要的未来研究领域，有赖于经验过程理论和统计估计理论的深入进展。这些进展带来了新的挑战，同时也是推动统计实践向前发展所必需的。

\section{注记与延伸阅读}

关于独立受试者层面数据的因果推断文献综述，可参阅该主题的若干部著作：Rubin (2006)、Pearl (2009a)、van der Laan and Robins (2003)、Tsiatis (2006)，以及 van der Laan and Rose (2011)。关于个体间存在因果干扰时不同类型因果效应的定义，参见 Halloran and Struchiner (1995)、Hudgens and Halloran (2008)、VanderWeele et al. (2012b)，以及 Tchetgen Tchetgen and VanderWeele (2012)。由于缺乏一般的方法学框架，过去许多研究为了简便，假设不存在个体间干扰。Sobel (2006) 通过实际情形展示了这一假设的风险，说明忽略个体间干扰，可能对干预的有效性得出错误结论。关于整群随机试验和群层面观察性研究的综述，还可参见 Donner and Klar (2000)、Hayes and Moulton (2009)，以及 Campbell et al. (2007)。

配套技术报告 van der Laan (2012) 包含多项补充结果，例如弱化序贯条件独立性假设：仍严格限制相依程度，但允许在给定观测过去之后，一个受试者仍可依赖至多 $K$ 个其他受试者；以及只观察完整因果相连总体的一个随机样本等。这一框架更一般的扩展，将是我们未来的研究目标。弱化基线协变量 $\mathbf W$ 独立性假设的另一种可能方法，是直接将目标参数定义为给定 $\mathbf W$ 的条件参数。van der Laan (2014a) 对这个估计量作了详细分析。

我们为可能存在网络相依的个体构建的半参数模型，推广了因果推断文献中针对独立个体的模型。尽管本章的因果模型涉及单组个体，但它涵盖可将个体划分为多个因果独立子组的情形。此外，我们的模型也包含分组序贯自适应设计：试验中某个个体的处理分配，取决于先前入组个体的已观测信息（Hu and Rosenberger 2006；van der Laan 2008a；Chambaz and van der Laan 2011a,b）。我们的模型还允许一个个体的结局，是其他个体所接受处理的函数。后者在因果推断文献中称为个体间干扰。因此，本文提出的因果模型，既推广了现有针对独立个体的因果模型，也推广了包含以往研究过的个体间因果依赖的因果模型。最后，我们指出，这些模型及相应方法，还可用于建立一种方法，评估针对网络的干预对各个体结局平均值的因果效应。例如，可能希望了解：通过某种干预改变社区的网络结构，如增加社区中特定个体之间的连接，会使社区层面的结局如何变化。此时，必须将处理节点定义为集合 $F_i$ 的属性，使处理的改变对应于网络结构的改变。最后，也可能关注既改变网络、又向部分个体分配暴露的干预的因果效应。

Aronow and Samii (2013) 给出了一种在一般个体间干扰下估计平均因果效应的 IPTW 方法，依靠实验设计生成所需的广义倾向评分（generalized propensity scores）。在已知广义倾向评分的前提下，作者还针对该 IPTW 估计量的真实条件方差，给出了有限样本下具有正偏差的估计量；这一真实方差不可识别，条件化所针对的是底层反事实。最后，作者考虑了观察多个来自子总体的独立样本时的渐近性质：子总体个数趋于无穷，每个样本均允许存在其一般形式的个体间干扰。他们的创新方法依赖于定义一个暴露模型，将 $N$ 个个体的处理节点，以及所指定的个体 $i$ 的特征，映射为个体 $i$ 的广义暴露（generalized exposure）。例如，除个体 $i$ 自身的暴露外，使用者还可将广义暴露定义为个体 $i$ 的朋友的暴露向量。对每个个体 $i$，该方法定义与某种静态干预对应的反事实结局：将这种广义暴露设为一个特定值，且对每个 $i$ 都设置为同一值；随后，将反事实平均结局定义为这些个体特定反事实的平均值的期望。它使用广义暴露条件概率的倒数加权，从而得到这些反事实结局平均值之期望的无偏估计量。

我们的模型将观察许多相互独立的个体簇的情形包含为特例。不过，通过采用更一般的条件独立性假设，我们还允许仅观察一个由相互连接个体构成的总体时，进行渐近统计推断。此外，我们用针对 $N$ 个个体特定暴露的随机干预来定义因果量，允许比个体间干扰更一般的依赖关系，并开发了与上述 IPTW 类型估计量迥异的高度有效估计量。总体而言，这些特点使我们的方法区别于 Aronow and Samii (2013)。

\medskip
\noindent\textbf{致谢} 本研究得到美国国立卫生研究院（NIH）项目 R01 AI074345-05 的资助。
